\chapter{The Geometry of Functions: Signals as Vectors}

\chapterline

Chapter~2 laid out the continuous and discrete representations of time and built signals out of impulses.
This chapter treats a whole signal as a single object. The rest of the book depends on
one geometric idea: a signal -- a sound, an
image, a stream of acceleration readings from a sensor on your wrist -- is a \emph{vector}, a genuine element of a vector space rather than a
metaphorical one, with all
the machinery of arrows, lengths, angles, and projections that you already know from two and three
dimensions. Seen this way, the Fourier transform is no longer a frightening integral; it
measures the components of a signal along a particular
set of perpendicular directions. The directions happen to be pure tones.

This chapter builds that picture from the ground up. We will start with arrows in the plane, watch
them grow extra coordinates until they become signals, and discover that the same three
operations -- adding, scaling, and taking inner products -- carry over completely. Along the way we
will meet the concepts that organize the rest of the book: \emph{energy} (a length),
\emph{correlation} (an angle), \emph{orthogonality} (perpendicularity), and \emph{orthonormal bases}
(perpendicular rulers you measure with). At the end, we will see that the sinusoids are the most
useful ruler of all, which leads directly to Fourier analysis.

\begin{focusbox}
The single sentence to carry out of this chapter: \textbf{analyzing a signal means projecting it
onto a basis, and a Fourier transform is exactly that, where the basis vectors are sinusoids.}
Everything else -- norms, inner products, orthogonality -- is the apparatus that makes ``project
onto'' precise.
\end{focusbox}

\section{From arrows to signals}
\label{ch03:arrows}

Start with the most familiar vector of all: an arrow in the plane, written as a pair of numbers
$\vec{v} = (v_1, v_2)$. You know three things you can do with it.

\begin{enumerate}
  \item \textbf{Add} two vectors: $(a_1,a_2) + (b_1,b_2) = (a_1+b_1,\, a_2+b_2)$. Geometrically,
    lay them tip-to-tail.
  \item \textbf{Scale} a vector by a number: $c\,(v_1,v_2) = (c v_1,\, c v_2)$. This stretches or
    flips the arrow without changing the line it lies on.
  \item \textbf{Measure} with the dot product: $\vec{a}\cdot\vec{b} = a_1 b_1 + a_2 b_2$. This one
    number encodes both lengths and the angle between them, through
    $\vec{a}\cdot\vec{b} = \norm{\vec{a}}\,\norm{\vec{b}}\cos\theta$.
\end{enumerate}

Now add a third coordinate, getting points in
$\R^3$; nothing breaks -- adding, scaling, and the dot product $a_1b_1+a_2b_2+a_3b_3$ all still
work. Add a fourth, a tenth, a thousandth. The pictures in your head get harder to draw, but the
\emph{algebra} never changes. A vector in $\R^N$ is just a list of $N$ numbers, and the three
operations are component-by-component sums, component-by-component scalings, and a dot product that
sums $N$ products.

This is the step that connects linear algebra to signal processing. A discrete signal sampled at
$N$ instants \emph{is} a list of $N$ numbers. The reading you take from a wrist accelerometer at
times $t_0, t_1, \dots, t_{N-1}$ is the vector
\[
  \vec{x} = \big(x[0],\, x[1],\, \dots,\, x[N-1]\big) \in \R^N .
\]
A one-second snippet of audio at $44{,}100$ samples per second is a single vector living in
$\R^{44100}$. The index that we usually call ``time'' is playing exactly the role that the index
$1,2,3$ played for the arrow's coordinates. There is no difference in kind, only in dimension.

\begin{intuitionbox}
Instead of a wiggly curve, think of a waveform as \emph{one point in a very
high-dimensional space}. The whole one-second audio clip is a single dot somewhere in $\R^{44100}$.
Two clips that ``sound similar'' are two dots that are close together. Silence is the origin. Making
a sound twice as loud is scaling its vector. Mixing two sounds is vector addition. The geometry you
will spend this book learning is the geometry of that space.
\end{intuitionbox}

\begin{figure}[H]
\centering
\begin{tikzpicture}[scale=1.0]
  % left: an arrow in R^2
  \begin{scope}[xshift=0cm]
    \draw[->,gray!60] (-0.3,0) -- (3.2,0) node[right]{$v_1$};
    \draw[->,gray!60] (0,-0.3) -- (0,2.8) node[above]{$v_2$};
    \draw[->,sigblue,very thick] (0,0) -- (2.4,2.0) node[above left,sigblue]{$\vec{v}$};
    \node[anchor=north] at (1.4,-0.45) {a vector in $\R^2$};
  \end{scope}
  % right: a signal as a stack of coordinates
  \begin{scope}[xshift=6.0cm]
    \begin{axis}[width=7.2cm,height=4.3cm,
        xlabel={sample index $n$}, ylabel={$x[n]$},
        xmin=-0.5,xmax=12.5,ymin=-1.4,ymax=1.4,
        axis lines=box, xtick={0,4,8,12}, ytick={-1,0,1},
        tick label style={font=\footnotesize}, label style={font=\footnotesize}]
      \addplot+[ycomb,sigblue,mark=*,mark size=1.2pt,thick] coordinates {
        (0,0.0)(1,0.5)(2,0.87)(3,1.0)(4,0.87)(5,0.5)(6,0.0)
        (7,-0.5)(8,-0.87)(9,-1.0)(10,-0.87)(11,-0.5)(12,0.0)};
    \end{axis}
    \node[anchor=north] at (3.4,-1.0) {the same idea with 13 coordinates};
  \end{scope}
\end{tikzpicture}
\caption{Left: a vector in $\R^2$ has two coordinates. Right: a length-$13$ signal has thirteen
coordinates, $x[0],\dots,x[12]$. A signal is just a vector with a lot of components; the index $n$
plays the role of the coordinate label.}
\label{ch03:fig-arrow-to-signal}
\end{figure}

\begin{connectionbox}
This chapter is the bridge between the companion \emph{Linear Algebra From the Ground Up} and all
the frequency-domain ideas to come. If you have read that book, every term here -- vector space,
inner product, norm, orthonormal basis, projection -- is the same object you already met, now with
the components indexed by \emph{time} instead of by $1,2,3$. In particular, a central theorem
of finite-dimensional linear algebra (every vector is a unique combination of an orthonormal basis)
becomes, almost word for word, the Fourier representation of a signal.
\end{connectionbox}

\section{What makes a vector space (and why signals form one)}
\label{ch03:vector-space}

We have been waving our hands at ``vector space.'' Let us pin it down, because the abstraction is
exactly what lets the same theorems serve arrows and audio at once.

\begin{definition}[Vector space, informally]
\label{ch03:def-vs}
A (real) vector space is a set of objects you can \emph{add} and \emph{scale by real numbers}, such
that the obvious rules hold: addition is commutative and associative, there is a zero element, every
element has a negative, scaling distributes over addition, and $1\cdot\vec{x}=\vec{x}$. Any set with
these two operations behaving sensibly is a vector space, and its elements are called vectors --
whatever they happen to look like.
\end{definition}

The phrase ``whatever they happen to look like'' matters: the definition never mentions
arrows. So let us check signals against it. Take the set of all real sequences of length $N$. Adding
two of them samplewise gives another length-$N$ sequence. Scaling one by a constant gives another.
The all-zeros sequence is the zero vector. Negatives, associativity, distributivity -- all hold,
because they hold component by component. So:

\begin{proposition}
\label{ch03:prop-RN}
The set of length-$N$ real signals, with samplewise addition and scalar multiplication, is the
vector space $\R^N$.
\end{proposition}

Nothing here required $N$ to be finite in spirit. The set of all infinite sequences
$\{x[n]\}_{n\in\Z}$, and even the set of all functions $x(t)$ of a continuous variable, satisfy
Definition~\ref{ch03:def-vs} just as well: you can add two functions, $(x+y)(t) = x(t)+y(t)$, and
scale a function, $(cx)(t)=c\,x(t)$. Continuous-time signals form a vector space too -- an
infinite-dimensional one. We will mostly reason in $\R^N$ (or its complex cousin $\C^N$) because the
hardware we build will always sample, and a finite list of numbers is what a microcontroller
actually holds. But keep in mind that the continuous case is the same animal with infinitely many
coordinates.

\begin{whybox}
Why insist on the abstract definition instead of just saying ``a signal is a list of numbers''?
Because the theorems we want -- the Pythagorean theorem, the existence of orthonormal bases, the
fact that projection minimizes error -- are proved \emph{only} from the vector-space axioms plus an
inner product. By verifying that signals satisfy the axioms, we inherit every one of those theorems
for free, without proving them a second time.
\end{whybox}

\section{The inner product: one number for length and angle}
\label{ch03:inner}

The dot product is central to plane geometry because it encoded both length and angle. We
need its generalization to signals. For two real length-$N$ signals $\vec{x}$ and $\vec{y}$, define

\begin{definition}[Inner product of signals]
\label{ch03:def-inner}
\[
  \inner{\vec{x}}{\vec{y}} \;=\; \sum_{n=0}^{N-1} x[n]\, y[n].
\]
For complex signals (which we will need the moment sinusoids become complex exponentials) we
conjugate the second argument:
\[
  \inner{\vec{x}}{\vec{y}} \;=\; \sum_{n=0}^{N-1} x[n]\, \overline{y[n]} ,
\]
where the overbar denotes the complex conjugate: if $z = a + \jj b$, then $\overline{z} = a - \jj b$.
Conjugation reflects a complex number across the real horizontal axis, negating its imaginary part
while preserving its real part. This ensures that the self-inner product of any complex number with itself
is its strictly real squared magnitude: $z\,\overline{z} = (a+\jj b)(a-\jj b) = a^2 + b^2 = \abs{z}^2$.
\end{definition}

This is the dot product, summed over time instead of over three spatial axes. The conjugate in the
complex case looks like a technicality but it is essential: it guarantees that
$\inner{\vec{x}}{\vec{x}}$ is a real, nonnegative number (a length squared had better not be
imaginary). Two familiar quantities follow from this definition.

\paragraph{It gives a length (energy).} Setting $\vec{y}=\vec{x}$,
\[
  \inner{\vec{x}}{\vec{x}} = \sum_{n} \abs{x[n]}^2 \;=\; \norm{\vec{x}}^2 .
\]
The right side is the \emph{energy} of the signal, and its square root is the norm, the length of the
vector. A loud sound is a long vector; a quiet sound is a short one. The match is literal: the physical
energy delivered by a voltage $x(t)$ into a unit resistor really is $\int x(t)^2\,\dd t$, the
continuous analog of $\sum x[n]^2$.

\paragraph{It gives an angle (correlation).} Just as in the plane,
\[
  \inner{\vec{x}}{\vec{y}} = \norm{\vec{x}}\,\norm{\vec{y}}\,\cos\theta,
  \qquad\text{so}\qquad
  \cos\theta = \frac{\inner{\vec{x}}{\vec{y}}}{\norm{\vec{x}}\,\norm{\vec{y}}}.
\]
The cosine of the ``angle between two signals'' is exactly what statisticians call the (normalized)
\emph{correlation} between them. Two signals that point the same way ($\theta=0$) are perfectly
correlated; two that point oppositely are anti-correlated; and two at right angles
($\inner{\vec{x}}{\vec{y}}=0$) are \emph{uncorrelated} -- they share nothing.

\begin{happeningbox}
When you ``compute the correlation'' of a candidate signal with a template -- say, matching a
recorded swipe gesture against a stored prototype -- you are literally measuring the cosine of the
angle between two vectors in a several-hundred-dimensional space. A big positive inner product means
the angle is small means ``these point the same way'' means ``this looks like that.'' Correlation
detectors, matched filters, and (soon) the Fourier transform are all the same geometric act: take an
inner product and see how much one vector leans along another.
\end{happeningbox}

\begin{example}[Inner products of short signals]
\label{ch03:ex-inner}
Let $\vec{x}=(1,2,3,4)$ and $\vec{y}=(4,3,2,1)$, both in $\R^4$.
\[
  \inner{\vec{x}}{\vec{y}} = 1\cdot4 + 2\cdot3 + 3\cdot2 + 4\cdot1 = 20,
\]
\[
  \norm{\vec{x}} = \sqrt{1+4+9+16}=\sqrt{30},\quad \norm{\vec{y}}=\sqrt{30},
  \quad \cos\theta = \frac{20}{30}\approx 0.667,
\]
so $\theta\approx 48^\circ$. The two signals are similar but not identical -- one rises, the other
falls, yet both have the same energy and a substantial overlap. Now try $\vec{u}=(1,-1,1,-1)$ and
$\vec{w}=(1,1,1,1)$: $\inner{\vec{u}}{\vec{w}} = 1-1+1-1 = 0$. These two are \emph{orthogonal}: a
fast alternation and a constant share nothing. We will return to this example; it is a tiny Fourier basis.
\end{example}

\begin{trapbox}
The inner product is \emph{not} elementwise multiplication, and it is \emph{not} the product of the
two signals. It is the single number you get after multiplying samplewise \emph{and then summing}.
Forgetting the sum (leaving a whole signal where a scalar should be) is the most common beginner
slip, and it makes ``project onto'' meaningless. Inner product in, one number out.
\end{trapbox}

\section{Norm, energy, and distance}
\label{ch03:norm}

Because the norm came from an inner product, it behaves like a true length and gives us a notion
of \emph{distance} between signals:
\[
  d(\vec{x},\vec{y}) = \norm{\vec{x}-\vec{y}} = \sqrt{\sum_n \abs{x[n]-y[n]}^2}.
\]
This is just the Euclidean (Pythagorean) distance, now in $N$ dimensions. It answers ``how different
are these two signals?'' with a single number, and it is the quantity that least-squares denoisers,
compressors, and regression models work to minimize. When we say a reconstructed signal is
``close to'' the original, we mean this distance is tiny.

One immediate payoff of the inner product is the Pythagorean theorem, and it is worth stating
because the idea of spectral energy depends on it.

\begin{theorem}[Pythagoras for signals]
\label{ch03:thm-pythagoras}
If $\vec{x}$ and $\vec{y}$ are orthogonal, $\inner{\vec{x}}{\vec{y}}=0$, then
\[
  \norm{\vec{x}+\vec{y}}^2 = \norm{\vec{x}}^2 + \norm{\vec{y}}^2 .
\]
\end{theorem}

\begin{proof}
Expand using bilinearity of the inner product:
$\norm{\vec{x}+\vec{y}}^2 = \inner{\vec{x}+\vec{y}}{\vec{x}+\vec{y}}
= \norm{\vec{x}}^2 + 2\inner{\vec{x}}{\vec{y}} + \norm{\vec{y}}^2$, and the middle term vanishes.
\end{proof}

\begin{aestheticbox}
Theorem~\ref{ch03:thm-pythagoras} is why we can speak of ``the energy in the 440~Hz component'' as a
separate, additive thing. If a signal decomposes into orthogonal pieces -- and sinusoids of
different frequencies, we will see, are orthogonal -- then the total energy is the \emph{sum} of the
energies of the pieces. No cross terms, no interference in the bookkeeping. This is an early form of Parseval's
theorem, and it is the reason a spectrum can be read straight off as ``how much energy lives at
each frequency.''
\end{aestheticbox}

\section{Orthogonality and orthonormal bases}
\label{ch03:ortho}

In the plane, the directions $\vec{e}_1=(1,0)$ and $\vec{e}_2=(0,1)$ are special: they are
perpendicular and unit length, and \emph{any} vector is a combination of them whose coefficients you
read off directly by inspection, $\vec{v}=v_1\vec{e}_1+v_2\vec{e}_2$. The coefficient $v_1$ is simply the
projection of $\vec{v}$ onto $\vec{e}_1$, i.e.\ $\inner{\vec{v}}{\vec{e}_1}$. We want this same
convenience for signals, with a basis we get to choose.

\begin{definition}[Orthonormal set]
\label{ch03:def-onb}
A set of vectors $\{\vec{b}_0,\vec{b}_1,\dots,\vec{b}_{N-1}\}$ is \emph{orthonormal} if
\[
  \inner{\vec{b}_k}{\vec{b}_m} = \delta_{km} =
  \begin{cases} 1 & k=m \\ 0 & k\neq m, \end{cases}
\]
that is, each vector has unit length and any two distinct vectors are orthogonal. If the set also
spans the whole space (here, $N$ such vectors in an $N$-dimensional space), it is an
\emph{orthonormal basis}.
\end{definition}

The key advantage of an orthonormal basis is the \emph{analysis formula}. Suppose
$\vec{x}=\sum_m c_m \vec{b}_m$ for some unknown coefficients $c_m$. Take the inner product of both
sides with one particular basis vector $\vec{b}_k$:
\[
  \inner{\vec{x}}{\vec{b}_k}
  = \Big\langle \sum_m c_m \vec{b}_m,\ \vec{b}_k \Big\rangle
  = \sum_m c_m \inner{\vec{b}_m}{\vec{b}_k}
  = \sum_m c_m\,\delta_{mk}
  = c_k .
\]
Every term in the sum collapses except the one that matches, because the basis is orthonormal. So
the coefficient is recovered by a single inner product:

\begin{theorem}[Analysis and synthesis]
\label{ch03:thm-analysis}
For an orthonormal basis $\{\vec{b}_k\}$, every signal $\vec{x}$ can be written
\[
  \underbrace{c_k = \inner{\vec{x}}{\vec{b}_k}}_{\text{analysis: extract a coefficient}},
  \qquad\qquad
  \underbrace{\vec{x} = \sum_{k=0}^{N-1} c_k\,\vec{b}_k}_{\text{synthesis: rebuild the signal}} .
\]
\end{theorem}

\begin{happeningbox}
To find ``how much of direction $\vec{b}_k$ is in my signal,'' you take
\emph{one inner product}. To rebuild the signal, you add up each direction weighted by how much of it
is present. \textbf{Analysis is a pile of inner products; synthesis is a weighted sum.} The Fourier
transform will be precisely this, with the $\vec{b}_k$ chosen to be sinusoids. The next several chapters fill in
the details.
\end{happeningbox}

\begin{whybox}
Why does orthonormality matter so much -- why not use any old basis? With a non-orthogonal basis you
\emph{can} still write $\vec{x}=\sum_m c_m \vec{b}_m$, but to find the coefficients you must solve a
coupled $N\times N$ linear system, because the inner-product trick above leaves cross terms
$\inner{\vec{b}_m}{\vec{b}_k}\neq 0$. Orthonormality \emph{diagonalizes} that system: each
coefficient is computed independently, by its own inner product, with no interference from the
others. That independence is what makes the Fast Fourier Transform fast and what makes a spectrum
interpretable channel by channel.
\end{whybox}

\subsection{A worked basis change}
\label{ch03:basis-change}

Let us make Theorem~\ref{ch03:thm-analysis} concrete in $\R^2$ before applying it to signals. Take
the rotated orthonormal basis
\[
  \vec{b}_0 = \tfrac{1}{\sqrt2}(1,1), \qquad \vec{b}_1 = \tfrac{1}{\sqrt2}(1,-1).
\]
Check: $\norm{\vec{b}_0}^2 = \tfrac12(1+1)=1$, likewise $\vec{b}_1$, and
$\inner{\vec{b}_0}{\vec{b}_1} = \tfrac12(1-1)=0$. Orthonormal. Now decompose $\vec{x}=(3,1)$:
\[
  c_0 = \inner{\vec{x}}{\vec{b}_0} = \tfrac{1}{\sqrt2}(3+1)=\tfrac{4}{\sqrt2}=2\sqrt2,\qquad
  c_1 = \inner{\vec{x}}{\vec{b}_1} = \tfrac{1}{\sqrt2}(3-1)=\sqrt2 .
\]
Reconstruct: $c_0\vec{b}_0 + c_1\vec{b}_1 = 2\sqrt2\cdot\tfrac{1}{\sqrt2}(1,1) +
\sqrt2\cdot\tfrac{1}{\sqrt2}(1,-1) = (2,2)+(1,-1)=(3,1)=\vec{x}$. It rebuilds perfectly. And notice
energy is preserved: $\norm{\vec{x}}^2 = 9+1 = 10$, while
$c_0^2 + c_1^2 = 8 + 2 = 10$. The coefficients carry exactly the same energy as the samples; we have
only rotated our point of view. (That is exactly the basis of Example~\ref{ch03:ex-inner}, scaled to
unit length -- the ``average'' direction $\vec b_0$ and the ``difference'' direction $\vec b_1$. It
is the smallest possible Fourier transform, and also the simplest wavelet.)

\begin{figure}[H]
\centering
\begin{tikzpicture}[scale=1.1]
  \draw[->,gray!55] (-0.4,0) -- (3.6,0) node[right]{$x_1$};
  \draw[->,gray!55] (0,-1.9) -- (0,2.4) node[above]{$x_2$};
  % standard basis grid hint
  \draw[gray!25,dashed] (3,0)--(3,1)--(0,1);
  % rotated basis directions
  \draw[->,structurepurple,very thick] (0,0) -- (1.6,1.6) node[above right]{$\vec{b}_0$};
  \draw[->,structurepurple,very thick] (0,0) -- (1.6,-1.6) node[below right]{$\vec{b}_1$};
  % the vector x
  \draw[->,sigblue,very thick] (0,0) -- (3,1) node[above right]{$\vec{x}=(3,1)$};
  % projection onto b0
  \draw[targetgreen,thick,dashed] (3,1) -- (2.0,2.0);
  \node[targetgreen,font=\footnotesize] at (2.9,1.95) {$c_0$ along $\vec{b}_0$};
\end{tikzpicture}
\caption{Changing basis is rotating your rulers. The same point $\vec{x}=(3,1)$ has coordinates
$(3,1)$ in the standard basis and $(2\sqrt2,\sqrt2)$ in the rotated orthonormal basis
$\{\vec{b}_0,\vec{b}_1\}$. The dashed green line is the projection that yields $c_0$. The point did
not move; only our description of it did.}
\label{ch03:fig-basis-change}
\end{figure}

\section{Projection: the geometry of ``how much of this is in that''}
\label{ch03:projection}

The coefficient $c_k=\inner{\vec{x}}{\vec{b}_k}$ is a \emph{projection}. It deserves its own
section because projection is also the right way to think about approximation, and approximation is
what real systems do when they cannot keep every coefficient.

The projection of $\vec{x}$ onto a unit vector $\vec{b}$ is the vector
$\inner{\vec{x}}{\vec{b}}\,\vec{b}$: it is the shadow $\vec{x}$ casts on the line through $\vec{b}$.
The leftover, $\vec{x} - \inner{\vec{x}}{\vec{b}}\,\vec{b}$, is orthogonal to $\vec{b}$ -- it is the
part of $\vec{x}$ that the direction $\vec{b}$ simply cannot see. This splits any signal into ``the
part along $\vec{b}$'' plus ``the part perpendicular to $\vec{b}$,'' and by
Theorem~\ref{ch03:thm-pythagoras} their energies add.

\begin{theorem}[Projection minimizes error]
\label{ch03:thm-bestapprox}
Among all scalar multiples $\alpha\vec{b}$ of a unit vector $\vec{b}$, the one closest to $\vec{x}$
(smallest $\norm{\vec{x}-\alpha\vec{b}}$) is $\alpha = \inner{\vec{x}}{\vec{b}}$. More generally, the
best approximation of $\vec{x}$ using only the directions $\{\vec{b}_k : k\in S\}$ is obtained by
keeping their projection coefficients $c_k=\inner{\vec{x}}{\vec{b}_k}$ and dropping the rest.
\end{theorem}

\begin{whybox}
This is the theoretical basis of lossy compression and of bandlimiting. When MP3 throws away
inaudible frequency components, or when a low-pass filter discards high-frequency directions, it is
\emph{projecting} the signal onto a smaller set of basis vectors. Theorem~\ref{ch03:thm-bestapprox}
guarantees this projection is the closest possible approximation within those kept directions -- you
cannot do better with the same budget of components. Keeping the largest coefficients and discarding
the smallest is provably the energy-optimal thing to do in an orthonormal basis.
\end{whybox}

\begin{example}[Approximating a ramp with a few basis vectors]
\label{ch03:ex-approx}
Take the ramp $\vec{x}=(1,2,3,4)$ from Example~\ref{ch03:ex-inner} and the orthonormal ``average''
direction $\vec{b}_0=\tfrac12(1,1,1,1)$ (length $1$, since $4\cdot(\tfrac12)^2=1$). Its coefficient
is $c_0=\inner{\vec{x}}{\vec{b}_0}=\tfrac12(1+2+3+4)=5$. The projection is
$c_0\vec{b}_0 = 5\cdot\tfrac12(1,1,1,1)=(2.5,2.5,2.5,2.5)$ -- the best constant approximation of the
ramp is its average, $2.5$, as expected. The residual $(-1.5,-0.5,0.5,1.5)$ is the
part of the ramp that a constant cannot represent, and it is orthogonal to $\vec{b}_0$ (its entries
sum to zero). To capture the residual we would need a basis vector that tilts -- a low ``frequency.''
\end{example}

\begin{figure}[H]
\centering
\begin{tikzpicture}
  % Left panel: vector projection geometry
  \begin{scope}[xshift=0cm]
    \draw[->,gray!50,thick] (-0.5,0) -- (4.8,0) node[right,font=\footnotesize]{$\operatorname{span}\{\vec{b}\}$};
    \draw[->,gray!30] (0,-0.5) -- (0,3.4);
    \draw[dashed,gray!60,thick] (-0.8,0) -- (4.6,0);
    \draw[->,structurepurple,very thick] (0,0) -- (1.5,0) node[below=2pt,font=\footnotesize]{$\vec{b}$};
    \draw[->,sigblue,line width=1.4pt] (0,0) -- (3.6,2.7) node[above right=1pt,font=\footnotesize]{$\vec{x}$};
    \draw[->,targetgreen,line width=1.4pt] (0,0) -- (3.6,0) node[below=4pt,font=\scriptsize]{$\hat{\vec{x}} = \inner{\vec{x}}{\vec{b}}\,\vec{b}$};
    \draw[->,specred,line width=1.2pt] (3.6,0) -- (3.6,2.7) node[midway,right=3pt,font=\footnotesize]{$\vec{e} = \vec{x} - \hat{\vec{x}}$};
    \draw[gray!70] (3.3,0) -- (3.3,0.3) -- (3.6,0.3);
    \node[anchor=north,font=\small] at (2.2,-0.9) {(a) Vector space: right triangle};
  \end{scope}

  % Right panel: waveform decomposition (Example 3.3)
  \begin{scope}[xshift=6.6cm]
    \begin{axis}[
      width=7.2cm, height=4.6cm,
      title={Signal decomposition $\vec{x} = \hat{\vec{x}} + \vec{e}$},
      title style={font=\small},
      xlabel={sample index $n$}, ylabel={amplitude},
      xmin=-0.5, xmax=3.5, ymin=-2.2, ymax=6.8,
      xtick={0,1,2,3}, ytick={-2,0,2,4},
      axis lines=box,
      tick label style={font=\footnotesize}, label style={font=\footnotesize},
      legend style={font=\scriptsize, at={(0.03,0.96)}, anchor=north west, fill opacity=0.9},
    ]
      \addplot[sigblue, thick, mark=*, mark size=1.5pt] coordinates {(0,1) (1,2) (2,3) (3,4)};
      \addlegendentry{original $\vec{x}$}
      \addplot[targetgreen, thick, dashed, mark=square*, mark size=1.3pt] coordinates {(0,2.5) (1,2.5) (2,2.5) (3,2.5)};
      \addlegendentry{projection $\hat{\vec{x}}$}
      \addplot[specred, thick, dotted, mark=triangle*, mark size=1.5pt] coordinates {(0,-1.5) (1,-0.5) (2,0.5) (3,1.5)};
      \addlegendentry{residual $\vec{e} \perp \vec{b}$}
    \end{axis}
    \node[anchor=north,font=\small] at (3.1,-0.8) {(b) Waveforms: average + residual};
  \end{scope}
\end{tikzpicture}
\caption{Orthogonal projection of a signal vector. (a) In the geometry of signal space, projecting $\vec{x}$ onto the unit direction $\vec{b}$ drops a perpendicular shadow $\hat{\vec{x}} = \inner{\vec{x}}{\vec{b}}\vec{b}$. The error vector $\vec{e} = \vec{x} - \hat{\vec{x}}$ is strictly perpendicular to $\vec{b}$, forming a right triangle with $\|\vec{x}\|^2 = \|\hat{\vec{x}}\|^2 + \|\vec{e}\|^2$. (b) The sample-by-sample waveforms from Example~\ref{ch03:ex-approx}: the best constant approximation $\hat{\vec{x}}$ captures the average level $2.5$, while the perpendicular residual $\vec{e}$ captures the ramp variation, summing to the original signal at every index.}
\label{ch03:fig-projection}
\end{figure}

\begin{walkbox}
Before reading on, predict: if you projected that same ramp onto a single one-cycle cosine instead of
the constant, would you get a large or small coefficient? (Hint: the ramp is mostly a steady rise
plus a gentle bow; a cosine over one window also rises then falls. They overlap, but not perfectly.)
We will return to this when the sinusoidal basis arrives -- the point is that \emph{each} basis
direction answers a different question about the same signal.
\end{walkbox}

\section{Sinusoids as an orthogonal basis: the start of Fourier analysis}
\label{ch03:sinusoid-basis}

We now choose our basis. Of all the orthonormal rulers we could lay down in signal space, one set is
especially useful: the \emph{sinusoids}. The most important reason (they are the eigenvectors of every
time-invariant system) is developed in Chapter~12, but one reason we can verify right now,
using only the inner product: \textbf{sinusoids of different frequencies are orthogonal.}

Consider the discrete signals built from a fundamental frequency that fits a whole number of cycles
into $N$ samples:
\[
  \vec{c}_k[n] = \cos\!\Big(\tfrac{2\pi k n}{N}\Big), \qquad
  \vec{s}_k[n] = \sin\!\Big(\tfrac{2\pi k n}{N}\Big), \qquad n=0,\dots,N-1 .
\]
The integer $k$ counts how many full cycles the sinusoid completes across the window. The claim is
that for $k\neq m$ (within range),
\[
  \inner{\vec{c}_k}{\vec{c}_m} = \sum_{n=0}^{N-1}
  \cos\!\Big(\tfrac{2\pi k n}{N}\Big)\cos\!\Big(\tfrac{2\pi m n}{N}\Big) = 0 ,
\]
and similarly cosines are orthogonal to sines, and distinct sines to each other. The cleanest way to
prove this -- and the reason complex exponentials are used throughout the book -- is to package
cosine and sine into a single rotating complex exponential
$\vec{e}_k[n]=e^{\jj 2\pi k n/N} = \cos(2\pi kn/N) + \jj\sin(2\pi kn/N)$. Then

\begin{theorem}[Orthogonality of the discrete complex exponentials]
\label{ch03:thm-dft-ortho}
For $k,m \in \{0,1,\dots,N-1\}$,
\[
  \inner{\vec{e}_k}{\vec{e}_m}
  = \sum_{n=0}^{N-1} e^{\jj 2\pi k n/N}\, \overline{e^{\jj 2\pi m n/N}}
  = \sum_{n=0}^{N-1} e^{\jj 2\pi (k-m) n/N}
  = \begin{cases} N & k=m,\\ 0 & k\neq m. \end{cases}
\]
\end{theorem}

\begin{proof}
If $k=m$ each term is $e^0=1$ and the sum is $N$. If $k\neq m$, write $r=e^{\jj 2\pi(k-m)/N}$, a
root of unity with $r\neq 1$. The sum is a finite geometric series
$\sum_{n=0}^{N-1} r^n = \dfrac{r^N-1}{r-1}$ (which follows from the algebraic identity
$(r-1)\sum_{n=0}^{N-1} r^n = r^N - 1$). Since $k-m$ is a nonzero integer,
$r^N = \big(e^{\jj 2\pi(k-m)/N}\big)^N = e^{\jj 2\pi(k-m)} = \cos(2\pi(k-m)) + \jj\sin(2\pi(k-m)) = 1$.
The numerator $r^N - 1$ is thus identically zero while the denominator $r-1 \neq 0$.
Hence the sum vanishes.
\end{proof}

\begin{happeningbox}
That geometric-series argument underlies all of Fourier analysis, and it has a
simple geometric meaning: summing $N$ equally spaced points around the unit circle gives the centroid of a
regular polygon, which sits at the center -- zero -- unless all the points coincide at $1$ (the case
$k=m$). The phasors cancel by symmetry. \emph{Different frequencies cancel; the matching frequency
reinforces.} That cancel-versus-reinforce is exactly how a Fourier transform ``picks out'' the
frequency present in a signal.
\end{happeningbox}

\begin{figure}[H]
\centering
\begin{tikzpicture}[scale=1.0]
  % left: k != m, 6 phasors evenly spaced -> cancel
  \begin{scope}[xshift=0cm]
    \draw[gray!40] (0,0) circle (1.3);
    \draw[->,gray!55] (-1.6,0) -- (1.6,0);
    \draw[->,gray!55] (0,-1.6) -- (0,1.6);
    \foreach \a in {0,60,120,180,240,300}{
      \draw[->,specred,thick] (0,0) -- (\a:1.3);
      \fill[specred] (\a:1.3) circle (1.2pt);
    }
    \fill[black] (0,0) circle (1.6pt);
    \node[anchor=north] at (0,-1.8) {$k\neq m$: sum $=0$};
  \end{scope}
  % right: k = m, all phasors at angle 0 -> add up
  \begin{scope}[xshift=5.2cm]
    \draw[gray!40] (0,0) circle (1.3);
    \draw[->,gray!55] (-1.6,0) -- (1.6,0);
    \draw[->,gray!55] (0,-1.6) -- (0,1.6);
    \draw[->,sigblue,very thick] (0,0) -- (1.3,0);
    \node[sigblue,font=\footnotesize] at (0.65,0.28) {$\times N$};
    \fill[sigblue] (1.3,0) circle (1.4pt);
    \node[anchor=north] at (0,-1.8) {$k=m$: sum $=N$};
  \end{scope}
\end{tikzpicture}
\caption{The geometric series of Theorem~\ref{ch03:thm-dft-ortho} as phasors on the unit circle.
When $k\neq m$ the terms $e^{\jj 2\pi(k-m)n/N}$ are equally spaced points around the circle (here
$N=6$); their average is the polygon's centroid, the origin, so they cancel. When $k=m$ every term
sits at angle $0$ and they reinforce to give $N$. Cancel versus reinforce is how the transform picks
out a frequency.}
\label{ch03:fig-phasor-sum}
\end{figure}

Figure~\ref{ch03:fig-phasor-sum} shows the two cases side by side.
So the $N$ complex exponentials $\{\vec{e}_0,\dots,\vec{e}_{N-1}\}$ are mutually orthogonal, each
with norm $\sqrt{N}$. Dividing by $\sqrt N$ makes them orthonormal:
$\vec{b}_k = \tfrac{1}{\sqrt N}\vec{e}_k$. By
Theorem~\ref{ch03:thm-analysis}, any signal expands as a sum of these, and the coefficients are inner
products. That expansion \emph{is} the Discrete Fourier Transform. We have, in effect, already
written it down:
\[
  X[k] = \inner{\vec{x}}{\vec{e}_k} = \sum_{n=0}^{N-1} x[n]\, e^{-\jj 2\pi k n/N},
  \qquad
  x[n] = \frac{1}{N}\sum_{k=0}^{N-1} X[k]\, e^{\jj 2\pi k n/N}.
\]
The analysis equation on the left extracts ``how much of frequency $k$ is in $\vec{x}$''; the
synthesis equation on the right rebuilds the signal from its frequency content. The factor $1/N$ is
the bookkeeping for our basis having norm $\sqrt N$ instead of $1$. Chapter~8 unpacks all of this,
but the geometry is already in place: \textbf{the DFT is a change of basis to the orthogonal
sinusoids, computed by projection.}

\begin{figure}[H]
\centering
\begin{tikzpicture}
  \begin{axis}[width=12cm,height=5.0cm,
      title={Three orthogonal basis sinusoids ($N=16$)},
      title style={font=\small},
      xlabel={sample index $n$}, ylabel={amplitude},
      xmin=0,xmax=16,ymin=-1.35,ymax=1.35,
      axis lines=box,
      legend style={font=\tiny, at={(0.5,-0.24)}, anchor=north, legend columns=3},
      samples=17, domain=0:16,
      tick label style={font=\footnotesize}, label style={font=\footnotesize}]
    \addplot[sigblue,thick,mark=*,mark size=1pt]   {cos(deg(2*pi*1*x/16))};
    \addlegendentry{$\vec{c}_1$ (1 cycle)}
    \addplot[specred,thick,mark=square*,mark size=1pt] {cos(deg(2*pi*3*x/16))};
    \addlegendentry{$\vec{c}_3$ (3 cycles)}
    \addplot[structurepurple,thick,dashed,mark=triangle*,mark size=1.2pt] {sin(deg(2*pi*1*x/16))};
    \addlegendentry{$\vec{s}_1$ (sine, 1 cycle)}
  \end{axis}
\end{tikzpicture}
\caption{Three members of the orthogonal sinusoidal basis for $N=16$: a one-cycle cosine, a
three-cycle cosine, and a one-cycle sine. Each pair has inner product zero -- summing the samplewise
products over the window gives exactly $0$. These mutually perpendicular ``directions'' are the rulers
the Fourier transform measures along.}
\label{ch03:fig-basis}
\end{figure}

\begin{historybox}
The idea of treating an entire function as a single geometric vector was developed
in the early twentieth century by \textbf{David Hilbert}, \textbf{Maurice Fréchet}, and
\textbf{Frigyes Riesz}. For two centuries following Newton and Leibniz, mathematicians conceived
of a function strictly as an analytic formula or a continuous curve drawn on paper. When
\textbf{Joseph Fourier} announced in 1807 that an arbitrary piecewise-smooth function could be
represented as an infinite sum of smooth sinusoids, the Paris Academy (including Lagrange and
Laplace) pushed back fiercely: how could an infinite sum of smooth, differentiable waves
ever equal a curve with sharp corners or jump discontinuities? Even Cauchy stumbled in 1821,
publishing an erroneous ``proof'' that the limit of any convergent sum of continuous functions
must itself be continuous.

Resolving these questions required a new kind of geometry. In 1906, \textbf{Maurice Fréchet}
abstracted the notion of distance, creating metric spaces. Between 1906 and 1912, \textbf{David
Hilbert} investigated integral equations by treating sequences and functions as vectors with
infinitely many coordinates, showing that Euclidean concepts like orthogonality and projection
could be preserved if one defined the inner product as an integral $\int f(t)\overline{g(t)}\,\dd t$.
In 1907, \textbf{Frigyes Riesz} and \textbf{Ernst Sigismund Fischer} independently proved the
key result: the space $L^2$ of square-integrable functions is \emph{complete} under the
energy norm $\norm{f} = \sqrt{\int \abs{f}^2\,\dd t}$. An infinite sum of orthogonal basis
functions does not need to converge pointwise at every single instant; it converges in
\emph{energy}, meaning the total squared error vanishes. Riesz further established the
\emph{Riesz Representation Theorem}: every continuous linear measurement on a Hilbert space is
simply an inner product with a fixed signal. Their work made it possible to treat signal
analysis as Euclidean geometry in infinite dimensions.
\end{historybox}

\section{Function space and infinite dimensions}
\label{ch03:l2}

Everything above lived in $\R^N$ or $\C^N$ -- finitely many coordinates. What happens for a
continuous-time signal $x(t)$, which has a value at every one of uncountably many instants? The same
geometry survives if we make two natural substitutions: sums become integrals, and the index $n$
becomes the continuous variable $t$. The inner product of two signals on an interval $[a,b]$ is
\[
  \inner{x}{y} = \int_a^b x(t)\,\overline{y(t)}\,\dd t,
  \qquad \norm{x}^2 = \int_a^b \abs{x(t)}^2\,\dd t .
\]
A function is ``in the space'' precisely when this energy integral is finite. That space -- square-
integrable functions with the integral inner product -- is called $L^2$, and it is the natural home
of finite-energy signals. It is infinite-dimensional, but it keeps the two properties we cared most
about: an inner product (hence lengths, angles, orthogonality) and \emph{completeness} (Cauchy
sequences of signals converge to a signal in the space). A space with those properties is a
\emph{Hilbert space}, and $L^2$ is the prototype.

\begin{intuitionbox}
Think of $L^2$ as ``$\R^\infty$ done carefully.'' Each instant $t$ contributes a coordinate, so a
continuous signal is a point with infinitely many coordinates. The integral $\int x\,\overline y\,\dd
t$ is the natural limit of the finite dot product $\sum x[n]\overline{y[n]}$ as the samples get dense.
Almost every geometric statement you trust in $\R^N$ -- Pythagoras, projection minimizes error, an
orthonormal basis lets you read coefficients by inner products -- holds in $L^2$. The one new
subtlety is convergence: an infinite sum of basis vectors must be shown to actually settle down to
the signal, which is the technical content of ``Fourier series converge.''
\end{intuitionbox}

\begin{connectionbox}
In \emph{Calculus From the Ground Up} the integral was area under a curve and the limit of Riemann
sums. Here that same integral serves as an \emph{inner product}, the continuous limit of a dot
product: the expression you learned to compute as an area now measures the angle between two
signals. And in \emph{Probability From the Ground Up}, the expectation $\E[XY]$ of two zero-mean
random signals is \emph{also} an inner product -- which is why ``uncorrelated'' (probability) and
``orthogonal'' (geometry) are the same word. All three books are looking at one mathematical
structure from three angles.
\end{connectionbox}

\begin{formalbox}
One caution: not every function lives in $L^2$. A pure infinite sinusoid $\cos(2\pi f t)$
has infinite energy over all of $\R$ ($\int \cos^2 = \infty$), so strictly it is not an $L^2$
function -- which is exactly why the plain Fourier transform of a sinusoid is a (Dirac) spike rather
than an ordinary function, and why in practice we always analyze a \emph{windowed}, finite chunk that
\emph{is} square-integrable. The finite, sampled vectors a microcontroller actually stores never run
into this; they always live safely in $\C^N$. We flag the subtlety so the integrals later do not
surprise you, not because it will trouble the hardware.
\end{formalbox}

\section{A first numerical taste}
\label{ch03:numerics}

The whole chapter is computable in a few lines. The snippet below builds two short signals, measures
their inner product, energy, and the angle between them, then projects one signal onto the
constant direction -- exactly Example~\ref{ch03:ex-approx}. Run it and watch the geometry turn into
numbers.

\begin{lstlisting}[language=Python,caption={Signals as vectors: inner product, energy, angle, projection.},label={ch03:code-inner}]
import numpy as np

x = np.array([1., 2., 3., 4.])
y = np.array([4., 3., 2., 1.])

ip     = np.dot(x, y)                 # inner product  -> 20.0
energy = np.dot(x, x)                 # ||x||^2        -> 30.0
cos_th = ip / (np.linalg.norm(x) * np.linalg.norm(y))
print("inner product :", ip)
print("energy of x   :", energy)
print("cos(angle)    :", round(cos_th, 4), "-> angle deg:",
      round(np.degrees(np.arccos(cos_th)), 1))

# Project x onto the unit "average" direction b0 = (1,1,1,1)/2
b0 = np.ones(4) / 2.0                 # ||b0|| = 1
c0 = np.dot(x, b0)                    # coefficient    -> 5.0
proj = c0 * b0                        # best constant approx -> 2.5 each
residual = x - proj
print("c0, projection:", c0, proj)
print("residual.b0   :", round(np.dot(residual, b0), 12))  # ~0
\end{lstlisting}

The next snippet verifies Theorem~\ref{ch03:thm-dft-ortho} numerically: it forms the matrix of all
pairwise inner products of the complex-exponential basis and confirms it is $N$ times the identity --
the basis really is orthogonal, by computer as well as by hand.

\begin{lstlisting}[language=Python,caption={The sinusoidal basis is orthogonal: inner products vanish off the diagonal.},label={ch03:code-ortho}]
import numpy as np
N = 16
n = np.arange(N)
# basis vectors e_k[n] = exp(j 2 pi k n / N), stacked as columns
E = np.array([np.exp(1j*2*np.pi*k*n/N) for k in range(N)])   # E[k] is e_k
G = E.conj() @ E.T                                           # Gram matrix <e_k, e_m>
print(np.round(G.real, 6))   # = N on the diagonal, 0 elsewhere
print("max off-diagonal magnitude:", np.max(np.abs(G - N*np.eye(N))))
\end{lstlisting}

\begin{labbox}
Change $N$ to any value and rerun the second snippet; the off-diagonal magnitude stays at the level
of floating-point dust ($\sim 10^{-13}$), confirming exact orthogonality. Then replace the complex
exponentials with \emph{non}-harmonic frequencies (say $k+0.5$ cycles per window) and watch the Gram
matrix fill up with nonzero off-diagonal terms -- the orthogonality is special to frequencies that
complete a whole number of cycles in the window. This observation explains why ``spectral
leakage'' exists and why we will later apply windows.
\end{labbox}

\section{Toward the wrist tracker}
\label{ch03:wrist}

\begin{goalbox}
Here is how this chapter applies to the gesture controller. The IMU on your wrist streams
acceleration and angular velocity, sampled at, say, $200$~Hz. A half-second of one axis is a vector
in $\R^{100}$. Recognizing a gesture -- a flick, a twist, a tap -- comes down to two geometric
operations you now know:

\begin{itemize}
  \item \textbf{Correlation with templates.} Store a prototype vector for each gesture
    (Figure~\ref{ch03:fig-imu-traces}). A live window is classified by which template it has the
    largest (normalized) inner product with -- the smallest angle. That is a nearest-direction
    classifier, based directly on the cosine of Section~\ref{ch03:inner}.
  \item \textbf{Projection onto a frequency basis.} Many gestures differ in their \emph{rhythm} and
    \emph{tempo}, not their raw shape. Projecting the window onto the sinusoidal basis
    (Section~\ref{ch03:sinusoid-basis}) -- i.e.\ taking its DFT -- turns ``how fast does the wrist
    oscillate'' into a handful of coefficients that are far more stable to classify than the raw
    samples. The energy in a few frequency bins is a compact, rotation-of-viewpoint summary of the
    motion.
\end{itemize}

The same inner product does double duty: against time-domain templates it measures similarity of
shape; against sinusoids it measures frequency content. We will build both pipelines in later
chapters; this chapter gave you the vector space they live in.
\end{goalbox}

\begin{figure}[H]
\centering
\begin{tikzpicture}
  \begin{axis}[width=12cm,height=4.6cm,
      title={Accelerometer templates for gesture matching in $\R^{100}$},
      title style={font=\small},
      xlabel={sample index $n$ (half a second at $200$~Hz)}, ylabel={accel.\ (g)},
      xmin=0,xmax=100,ymin=-1.6,ymax=1.6,
      axis lines=box,
      legend style={font=\tiny, at={(0.5,-0.26)}, anchor=north, legend columns=2},
      samples=100, domain=0:100,
      tick label style={font=\footnotesize}, label style={font=\footnotesize}]
    % a "flick": a damped oscillation
    \addplot[sigblue,thick] {1.3*exp(-x/28)*sin(deg(2*pi*x/22))};
    \addlegendentry{a ``flick'' window}
    % a "double tap": two short bursts
    \addplot[specred,thick,dashed] {0.9*exp(-((x-20)^2)/30) - 0.9*exp(-((x-45)^2)/30)};
    \addlegendentry{a ``double tap'' window}
  \end{axis}
\end{tikzpicture}
\caption{Two half-second windows of a single accelerometer axis, each a vector in $\R^{100}$.
A flick is a brief damped oscillation; a double tap is a pair of sharp bursts. Classifying a live
window is choosing which stored template it has the smallest angle to -- a single normalized inner
product per template.}
\label{ch03:fig-imu-traces}
\end{figure}

\begin{buildbox}
A concrete starting sketch for an Arduino-class board reading an IMU: keep a circular buffer of the
last $N$ samples per axis, and every few milliseconds compute the inner product of that buffer with
each stored gesture template (a multiply-accumulate loop -- the one primitive every DSP chip is built
around). Normalize by the buffer's norm so loudness/intensity does not dominate. In fixed-point C the
core is:

\begin{lstlisting}[language=C,caption={Inner-product (matched-filter) core for an MCU.},label={ch03:code-mac}]
// returns sum_{i} buf[i]*tmpl[i]  (a single inner product / projection)
int32_t inner_product(const int16_t *buf, const int16_t *tmpl, int N) {
    int32_t acc = 0;
    for (int i = 0; i < N; ++i)
        acc += (int32_t)buf[i] * (int32_t)tmpl[i];   // multiply-accumulate
    return acc;
}
\end{lstlisting}

That loop is a direct implementation of Definition~\ref{ch03:def-inner}. Everything fancier -- the FFT, FIR
filtering, matched detection -- is built from millions of these.
\end{buildbox}

\begin{trapbox}
\textbf{What the wrist can actually sense.} A single wrist-worn IMU measures the motion of the
\emph{wrist}, full stop. No amount of geometry recovers individual finger joint angles from it: the
information simply is not in the signal. An IMU can reliably distinguish \emph{coarse, discrete}
gestures (a flick, a rotation, a double-tap, broad hand poses with characteristic motion), but
continuous finger pose is out of reach. Finger-level discrimination needs a different sensor that
sees muscle activity -- surface EMG or force myography (FMG) on the forearm -- and truly continuous
individual-finger tracking needs sensors \emph{on the fingers} (instrumented gloves or rings). We
will cover EMG and FMG in their own chapters. For now: signals-as-vectors is the right
framework, but choose a feature your sensor can actually observe, otherwise the inner product will only
measure noise.
\end{trapbox}

\section{Summary: Capabilities, Downsides, and Forward Roadmap}
\label{ch03:summary}

Putting the pieces of this chapter together: a signal is a vector; the inner product gives it a length
(energy), an angle to other signals (correlation), and a notion of perpendicularity (orthogonality);
an orthonormal basis lets you read off coordinates by single inner products (analysis) and rebuild the
signal by a weighted sum (synthesis); projection onto a chosen basis is the best approximation within
that basis; and the sinusoids form an orthogonal basis whose analysis-and-synthesis is the Fourier
transform. The rest of the book builds on these facts.

\begin{connectionbox}
\textbf{Capabilities, Downsides, and Forward Roadmap}
\begin{itemize}
  \item \textbf{What this made possible}: The entire apparatus of Euclidean geometry is now available for
  signals. Parseval's identity is the Pythagorean theorem; best approximation is orthogonal projection;
  and cross-correlation is measuring the angle between signal vectors. We no longer treat filtering
  or transforms as opaque equations, but as coordinate rotations and geometric projections.
  \item \textbf{Limitations}: Vector-space geometry gives you the machinery of
  bases and projections, but it has a major limitation: it does not tell you \emph{which} basis to
  choose. In $\R^N$ or $L^2$, there are infinitely many valid orthonormal bases (Walsh--Hadamard functions,
  Legendre polynomials, Haar wavelets, Chebyshev polynomials). Pure geometry cannot explain why
  sinusoids are central to signal processing. If you only look at signals as abstract
  vectors, all orthonormal bases look equally valid.
  \item \textbf{Where we go next}: To see why sinusoids dominate physical signals, we must study
  systems that operate on them. Chapter 4 examines sinusoids, phasors, and complex exponentials in
  depth. Then, in Chapters 11 and 12, we prove the central theorem of linear systems: sinusoids are the
  unique \emph{eigenfunctions} of Linear Time-Invariant (LTI) operators. Convolution in the time
  domain becomes simple pointwise multiplication in the Fourier basis. That property is what makes
  the Fourier basis the natural choice for analyzing LTI systems.
\end{itemize}
\end{connectionbox}

\section{Exercises}
\label{ch03:exercises}

\begin{exercise}
Let $\vec{x}=(2,-1,0,3)$ and $\vec{y}=(1,1,1,1)$ in $\R^4$. Compute $\inner{\vec{x}}{\vec{y}}$,
$\norm{\vec{x}}$, $\norm{\vec{y}}$, and the angle $\theta$ between them. Is $\vec{x}$ closer to being
parallel or perpendicular to the constant signal $\vec{y}$? What does $\inner{\vec{x}}{\vec{y}}$ tell
you about the \emph{average} of $\vec{x}$?
\end{exercise}

\begin{exercise}
Show directly from Definition~\ref{ch03:def-inner} that the inner product is linear in its first
argument: $\inner{a\vec{u}+b\vec{v}}{\vec{w}} = a\inner{\vec{u}}{\vec{w}} + b\inner{\vec{v}}{\vec{w}}$
for real scalars $a,b$. Where in the proof of Theorem~\ref{ch03:thm-pythagoras} was this used?
\end{exercise}

\begin{exercise}
Verify that $\vec{u}=(1,-1,1,-1)$ and $\vec{w}=(1,1,1,1)$ are orthogonal, and that together with
$\vec{p}=(1,1,-1,-1)$ and $\vec{q}=(1,-1,-1,1)$ they form an orthogonal set of four vectors in
$\R^4$. Normalize them to an orthonormal basis. (This is the $4$-point Walsh--Hadamard basis -- a
``square-wave'' cousin of the Fourier basis.)
\end{exercise}

\begin{exercise}
Using the orthonormal basis from the previous exercise, find the four coefficients $c_k =
\inner{\vec{x}}{\vec{b}_k}$ for $\vec{x}=(4,2,2,0)$, then reconstruct $\vec{x}$ from them to check.
Confirm Parseval: $\norm{\vec{x}}^2 = \sum_k c_k^2$.
\end{exercise}

\begin{exercise}
Prove the Cauchy--Schwarz inequality $\abs{\inner{\vec{x}}{\vec{y}}} \le \norm{\vec{x}}\,
\norm{\vec{y}}$ for real signals, by considering $\norm{\vec{x}-t\vec{y}}^2 \ge 0$ as a quadratic in
$t$ and requiring its discriminant to be nonpositive. Why does this guarantee that the ``angle''
defined by $\cos\theta = \inner{\vec{x}}{\vec{y}}/(\norm{\vec{x}}\norm{\vec{y}})$ is always a real
angle (i.e.\ $\abs{\cos\theta}\le 1$)?
\end{exercise}

\begin{exercise}
For $N=4$, write out the four complex-exponential basis vectors $\vec{e}_k[n]=e^{\jj 2\pi kn/4}$ as
explicit $4$-tuples (their entries are among $\{1,\jj,-1,-\jj\}$). Verify by direct summation that
$\inner{\vec{e}_1}{\vec{e}_3}=0$ but $\inner{\vec{e}_1}{\vec{e}_1}=4$, matching
Theorem~\ref{ch03:thm-dft-ortho}.
\end{exercise}

\begin{exercise}
Take the ramp $\vec{x}=(1,2,3,4)$. Its best constant approximation was $(2.5,2.5,2.5,2.5)$
(Example~\ref{ch03:ex-approx}). Compute the energy of the residual and confirm that
$\norm{\vec{x}}^2 = \norm{\text{projection}}^2 + \norm{\text{residual}}^2$, an instance of
Theorem~\ref{ch03:thm-pythagoras}. What fraction of the signal's energy is captured by the constant
term alone?
\end{exercise}

\begin{exercise}
Using the geometric-series identity, evaluate $\sum_{n=0}^{N-1} e^{\jj 2\pi n/N}$ for $N=5$ in two
ways: (a) symbolically, showing it is exactly $0$; (b) by plotting the five phasors on the unit
circle and arguing geometrically that they sum to the center. Connect this to the ``different
frequencies cancel'' claim of Section~\ref{ch03:sinusoid-basis}.
\end{exercise}

\begin{exercise}
(Continuous case.) On the interval $[0,1]$, show that $\cos(2\pi t)$ and $\cos(4\pi t)$ are
orthogonal under the integral inner product $\int_0^1 x(t)y(t)\,\dd t$, and compute the norm of each.
Why is $\int_0^1 \cos^2(2\pi t)\,\dd t = \tfrac12$ rather than $1$, and what unit-norm version of the
cosine would you use as a basis vector?
\end{exercise}

\begin{exercise}
(Sensor limits.) Suppose you record one accelerometer axis during a ``wrist flick'' and during a
``finger snap.'' Explain, in terms of which information lives in the IMU signal vector, why a matched
filter might distinguish the flick from rest but cannot tell you \emph{which finger} snapped. What
additional sensor (and what kind of signal vector) would carry the missing information? Cite the
distinction between IMU, EMG/FMG, and on-finger sensing.
\end{exercise}

\begin{exercise}
(Computational.) Modify the second code listing to build the real sinusoidal basis (cosines and sines
for $k=0,\dots,N/2$) for $N=8$ and form its Gram matrix. Identify which pairs are orthogonal and
which basis vectors are missing or redundant (think about $k=0$ and $k=N/2$). How many real basis
vectors do you need to span all of $\R^8$, and does that match the count of independent DFT
coefficients?
\end{exercise}
